\scanpage{51}
\mspage{21}
Il y a de plus des ``conditions de circonstance'' :

\textcircled{c} L'application $\Sigma^{<}(1)\xrightarrow{\mathrm{source}}\Sigma^{<}(0)$ est \inserted{propre et} fibrante (``condition \inserted{d'}isotrivialité''), cela implique que \deleted{\unsure{si}} si $d\subset d'$ est une inclusion d'ens finis tot. ordonnés, tel que $d$ soit segment initial de $d'$, alors $\Sigma^{\unsure{*|<}}(d')-\Sigma^{<}(d)$ est fibrante.

\textcircled{d} \deleted{Le} $\Sigma^{<}(0)$ est \deleted{\unsure{le}} une variété à bord, \uline{et} \inserted{$\Sigma^{<}(1)\xrightarrow{\mathrm{but}}\Sigma^{<}(0)$ a comme image son bord.} (``condition non singulière'') donc aussi les $\Sigma^{<}(d)$ (car un sous-espace bordant d'une variété à bord est une variété à bord \deleted{\unsure{qui doit}} \unsure{déterminer le bord}, de $=$ dim que celui-ci ?? \unsure{\}} ou on exige \deleted{\unsure{de plus la condition}} \inserted{que les $\Sigma^{<}(d)$ soient des} variétés à bord\ldots)
\margin{Dans tous les cas, l'image de $\Sigma^{<}(1)\xrightarrow{\mathrm{but}}\Sigma^{<}(0)$ est une partie loc\textsuperscript{t} bordante\ldots}
\margin{il faudrait vérifier que le bord de $\Sigma^{<}(d)$ est l'image des $\Sigma^{<}(d')$ pour $d'\supset d$ cofinaux \unsure{à $d$}. \deleted{\unsure{??, ça}} Si c'est vrai, ça doit être délicat ! (d'après le cas où $\Sigma$ est au dessus de $I=\unsure{\{1,2,3\}|\{0,1,2\}}$)}
\margin{\deleted{\unsure{Grosso}} c'est peut être une conséquence de c) et de d)\ldots}

\textcircled{e} (Condition \deleted{\unsure{de relation de}} \inserted{mixte} équisingulière-lisse)
\deleted{\illegible{2} $\unsure{\Sigma}$} Le morphisme
\[
\Sigma^{<}(1)\xrightarrow{\mathrm{source}}\Sigma^{<}(0)
\]
(qui est propre et fibrant par c)) \deleted{\unsure{a encore des} \unsure{plong.} \unsure{et des variétés}} est \inserted{à fibres des} variétés à bord, et le fibré des bords est \deleted{à bord, Il \unsure{est visible que de prendre un} \unsure{produit}} \inserted{l'image} de $\Sigma^{<}(2)$ par $\Sigma^{<}(2)-\Sigma^{<}(1)$ associé à $\{0,2\}\longleftrightarrow\{0,1,2\}$
\mspage{50}
